\subsection{Improvements to the methods}\label{sec:improve}

% Rubric: "discuss or implement possible solutions to improve the methods".
% Each row: a failure we observed, its cause, the fix we implemented, and the measured effect.
% [Ky] Sources: K-Means++ row results/hotspot/ablation.csv (init, Ecoli); D3 gap = silhouette minus
% shuffled-data silhouette at K=2 and K=5, results/hotspot/fake_data_check.csv.

\begin{table}[h]\centering\small
\caption{Implemented improvements. Each fix targets a failure observed in Section~\ref{sec:comparison}.}
\label{tab:improve}
\resizebox{\textwidth}{!}{\begin{tabular}{lllll}
\toprule
\textbf{Failure} & \textbf{Cause} & \textbf{Fix} & \textbf{Data} & \textbf{Before $\to$ after} \\
\midrule
Clusters split at $180^\circ$ & flat lat/lon & surface distance & D2 & 0 $\to$ 92 of 92 pairs kept \\
Poor local optimum & random seeding & K-Means++ & D4 & ARI 0.485 $\to$ 0.561 (30 seeds) \\
Sparse belts lost & one global $\varepsilon$ & HDBSCAN & D2 & Americas kept 20\% $\to$ 62\% \\
One feature dominates & unequal scales & standardisation & D4 & ARI 0.497 $\to$ 0.731 \\
Clusters are storey bands & recorded steps & jitter within step & D8 & band clusters 107 of 124 $\to$ 0 of 23 \\
Curved clusters cut & convex assumption & spectral clustering & D1 & ARI 0.481 $\to$ 1.000 \\
Degenerate $k$ chosen & silhouette alone & constraints + fake-data check & D2, D3 & D2 null gap 0.015 $\to$ 0.144; D3 0.089 $\to$ 0.113 \\
\bottomrule
\end{tabular}}
\end{table}

% Sources: results/earthquakes_2023/ablation_distance.csv, case_density_regions.csv,
% kmeans_k_selection.csv; results/hdb_resale/hdbscan_grid.csv (single_storey_clusters, 9 settings
% per data version).
\paragraph{Surface distance (D2).} Haversine distance for density methods and 3D unit vectors for
K-means measure what matters, the distance along the Earth's surface. The fix is exact, and its only
cost is that $\varepsilon$ must be given in kilometres and the coordinates converted. It changes the
answer for every method, which makes it a correctness fix rather than a tuning choice.

\paragraph{HDBSCAN instead of DBSCAN (D2).} HDBSCAN replaces the single radius with a hierarchy of
densities and keeps the clusters that persist longest, so sparse and dense belts can both survive. In the
Americas it clusters 62\% of earthquakes against 20\%. The costs are real: dense cores are trimmed (78\%
against 90\% in the south-west Pacific), each fit takes about 20 times longer (1.08\,s against 0.05\,s),
and \texttt{min\_cluster\_size} is a count, so the result changes with the size of the dataset.

\paragraph{Jitter within the recording step (D8).} Moving each stepped value by up to half its step
spreads exact ties without moving any sale into another band. Across the 9 HDBSCAN settings, 107 of the
124 clusters found on recorded data are single storey bands; on jittered data, 0 of 23 are. The fix costs
nothing at the chosen setting (recorded and jittered agree at ARI 0.886), but it adds randomness, so the
seed must be fixed and reported. It also shows that D8 has fewer real density clusters than the recorded
data suggest.

\paragraph{Fake-data check for choosing $K$ (D2).} The silhouette alone picks $K=2$ on D2 (0.544, the
highest of any $K$), but shuffled coordinates score 0.529 at the same $K$: the split follows the
distribution of latitude and longitude on their own, not their combination. Requiring a clear gap over the
fake data rules out $K=2$ and $K=3$. From $K=4$ the gap is 0.126 to 0.149 and almost flat, so the final
choice, $K=6$, rests on agreement between seeds (ARI 0.993). The check costs extra clustering runs on fake
data for each $K$.
\habegin
\paragraph{Proposed, not implemented (D6).} \hatag{} The density methods fail on D6 because 13
standardised features leave little contrast in density. Projecting onto a few principal components
before DBSCAN or HDBSCAN is the natural next step.
\haend

\paragraph{Spectral clustering for curved clusters (D1).} % results/datasets/make_moons_results.csv
K-Means++ fails on D1 because it assigns each point to the nearest centroid, so cluster boundaries
are straight lines. Spectral clustering replaces Euclidean distance with connectivity in a
10-nearest-neighbour graph: points along one moon are linked by short hops, the two moons are
not, and the eigenvectors of the graph Laplacian separate them. ARI rises from 0.481 to 1.000. The
cost is that $k$ and the graph must still be chosen, and the affinity matrix and eigenvectors limit
how far the method scales (Section~\ref{sec:runtime}).
\paragraph{Proposed, not implemented (D2).} D2 could be scored rather than inspected: assigning each
earthquake to the nearest segment of a published plate-boundary model would give labels for ARI, and
show directly whether a cluster follows one boundary or joins several. OPTICS~\cite{ankerst1999}, which
orders points by reachability so that clusters at several densities can be read from one plot, would
show whether a single ordering explains both the dense and the sparse belts.
\todo{Lam, Ky}{other proposals, e.g.\ approximate spectral clustering for large $n$.}
